\subsubsection*{I. Khái niệm sóng âm}
\paragraph{Sóng âm}
\notebox{Sóng âm là \textit{sóng cơ} lan truyền trong môi trường rắn, lỏng, khí.
\begin{itemize}
\item Sóng âm trong môi trường lỏng, khí là \textit{sóng dọc}.
\item Sóng âm trong môi trường rắn là \textit{sóng dọc} hoặc \textit{sóng ngang}.
\end{itemize}
}
\paragraph{Nguồn âm}
\notebox{Nguồn âm là các vật dao động phát ra âm
\begin{itemize}
\item Tần số của âm phát ra bằng tần số dao động của nguồn âm.
\end{itemize}
}
\paragraph{Phân loại âm}
\begin{center}
\begin{tabular}{|c|c|c|c|}
\hline 
\textbf{Phân loại âm} & \textbf{Âm nghe được (âm thanh)} & \textbf{Hạ âm} & \textbf{Siêu âm} \\ 
\hline 
\textbf{Tần số} & 16 Hz đến 20.000 Hz & < 16 Hz & > 20.000 Hz \\ 
\hline 
\end{tabular} 
\end{center}
%\paragraph{Môi trường truyền âm - Vận tốc âm}
%\begin{itemize}
%\item \textbf{Môi trường truyền âm}
%\begin{itemize}
%\item Môi trường truyền âm có thể là môi trường chất rắn, lỏng, hoặc khí. 
%\item Sóng âm\textit{ không truyền được trong môi trường chân không}.
%\item Những vật liệu như bông, nhung, tấm xốp, tấm kính truyền âm kém.
%\end{itemize}
%\item \textbf{Tốc độ truyền âm}
%\begin{itemize}
%\item Sóng âm truyền trong mỗi môi trường với một tốc độ hoàn toàn xác định.
%\item Vận tốc truyền âm phụ thuộc vào\textbf{ tính đàn hồi},\textbf{ mật độ} và \textbf{nhiệt độ} của môi trường truyền âm.
%\begin{align*}
%\boder{v_{\text{rắn}}>v_{\text{lỏng}}>v_{\text{khí}}}
%\end{align*}
%\end{itemize}
%\end{itemize}
\subsubsection*{II. Các đặc trưng vật lí của âm thanh}
\setcounter{paragraph}{0}
\paragraph{Tần số âm}
Tần số âm cũng là tần số của \textbf{nguồn âm}.
\Note{\textit{Tần số âm} là đặc trưng vật lí quan trọng nhất của âm}
\paragraph{Cường độ âm. Mức cường độ âm}
\begin{itemize}
\item Cường độ âm \bth{I} tại một điểm là đại lượng đo bằng năng lượng mà sóng âm truyền qua một đơn vị diện tích đặt vuông góc với phương truyền âm trong một đơn vị thời gian. 
\begin{align*}
\boder{I=\dfrac{W}{S.t}=\dfrac{P}{S}}\quad (W/m^2)
\end{align*}
\item Mức cường độ âm \bth{L} là đại lượng đo bằng loga thập phân của tỉ số giữa cường độ âm \bth{I} tại điểm đang xét và cường độ âm chuẩn \bth{I_0} (\bth{I_0=10^{-12}\ W/m^2} ứng với tần số \bth{f=1000\ Hz}). 
\begin{align*}
\boder{L=\lg\dfrac{I}{I_0}}\quad (\mathrm{B})
\end{align*}
\end{itemize}
\Note{\textit{Mức cường độ âm} là đặc trưng vật lí thứ hai của âm}
\paragraph{Âm cơ bản - Hoạ âm - Đồ thị dao động âm}
\begin{itemize}
\item Khi các nhạc cụ phát ra âm thì chúng đồng thời phát ra âm cơ bản có tần số $f_0$ và các hoạ âm có tần số $2f_0,\ 3f_0,\ 4f_0...$
\item Tập hợp âm cơ bản và hoạ âm tạo ra đồ thị dao động riêng của nhạc âm đó. 
\end{itemize}
\Note{\textit{Đồ thị dao động} là đặc trưng vật lí thứ ba của âm, phụ thuộc tần số và biên độ của các họa âm}
\subsubsection*{III. Các đặc trưng sinh lí của âm}
\setcounter{paragraph}{0}
\paragraph{Độ cao}
\textbf{Độ cao} là một đặc trưng sinh lí của âm, gắn liền với tần số âm.
\notebox{
\begin{itemize}
\item Âm cao (thanh) là âm có tần số âm lớn hơn.
\item Âm thấp (trầm) là âm có tần số âm nhỏ hơn. 
\end{itemize}}
\paragraph{Độ to}
\textbf{Độ to} là một đặc trưng sinh lí của âm, gắn liền với đặc trưng vật lí mức cường độ âm L.
\notebox{
\begin{itemize}
\item Độ to phụ thuộc cả vào \textit{tần số âm}.
\item Độ to của âm tăng theo mức cường độ âm nhưng không tỉ lệ thuận. 
\end{itemize}}
\paragraph{Âm sắc}
\textbf{Âm sắc} là đặc trưng sinh lí của âm giúp ta phân biệt được các âm phát ra từ các nguồn khác nhau. Âm sắc liên quan mật thiết đến đồ thị dao động âm.

%\subsubsection*{IV. Ngưỡng nghe. Ngưỡng đau. Miền nghe được}
%\begin{itemize}
%\item Giá trị cường độ âm \bth{I} bé nhất mà tai người còn cảm nhận được gọi là ngưỡng nghe. 
%\Note{
%Giá trị của ngưỡng nghe phụ thuộc vào từng tai người, cường độ âm và tần số âm.
%}
%\item Giá trị \bth{I} nào đó đủ lớn làm tai nghe có cảm giác nhức nhối, đau đớn thì gọi là ngưỡng đau. 
%\Note{
%\begin{itemize}
%\item Giá trị của ngưỡng đau không phụ thuộc vào tần số.
%\item Khi âm vượt qua ngưỡng đau thì tai ta vẫn nghe được âm.
%\end{itemize}
%}
%\item Miền \bth{I} nằm trong khoảng ngưỡng nghe và ngưỡng đau gọi là miền nghe được.
%\begin{align*}
%I_{\min}\le I\le I_{\max}
%\end{align*}
%\end{itemize}